% Source: https://artofproblemsolving.com/community/c3h1934648_geometry_problem
% Stable identity: aops-topic-1934648-q01
Scale so that $AC=1$; then $CD=DE=EB=1$ and $CB=3$.

\begin{asy}
usepackage("amsmath");
size(7cm);
defaultpen(fontsize(10pt));
pair C=(0,0),D=(1,0),Ep=(2,0),B=(3,0),A=(0,1);
fill(A--D--Ep--cycle,paleblue);
draw(A--C--B--cycle,linewidth(0.9)); draw(A--D); draw(A--Ep);
draw((0,0.12)--(0.12,0.12)--(0.12,0));
draw(arc(D,0.22,90+45,180),red); draw(arc(Ep,0.3,180-26.565,180),deepgreen); draw(arc(B,0.45,180-18.435,180),blue);
dot(A);dot(B);dot(C);dot(D);dot(Ep);
label("$A$",A,N);label("$B$",B,SE);label("$C$",C,SW);label("$D$",D,S);label("$E$",Ep,S);
label("$1$",(0,0.5),W);label("$1$",(0.5,0),S);label("$1$",(1.5,0),S);label("$1$",(2.5,0),S);
label("\mbox{}",(min(currentpicture,user=true)+max(currentpicture,user=true))/2,black); // keeps the text after the figure black when LaTeX typesets the labels inline
\end{asy}

\textbf{$\angle ADC=45^\circ$.} Triangle $ACD$ is a right isosceles triangle ($AC=CD=1$), so $\angle ADC=45^\circ$ and $AD=\sqrt2$.

\textbf{Similar triangles.} Compare triangles $DAE$ and $DBA$ (shaded triangle $ADE$ and the larger triangle $ADB$). They share the angle at $D$, since $E$ and $B$ lie on the same ray from $D$, and
\[\dfrac{DE}{DA}=\dfrac{1}{\sqrt2},\qquad \dfrac{DA}{DB}=\dfrac{\sqrt2}{2}=\dfrac{1}{\sqrt2}.\]
So $\triangle DEA\sim\triangle DAB$ (SAS), with $E\leftrightarrow A$ and $A\leftrightarrow B$. Hence $\angle DAE=\angle DBA=\angle ABC$.

\textbf{Exterior angle.} Since $D$ lies between $C$ and $E$, the angle $ADC$ is an exterior angle of triangle $ADE$, so
\[\angle ADC=\angle DAE+\angle AED=\angle ABC+\angle AEC .\]
Therefore
\[\angle ABC+\angle AEC+\angle ADC=2\angle ADC=2\cdot45^\circ=90^\circ.\ \blacksquare\]
(In terms of tangents: $\tan\angle ABC=\dfrac{1}{3}$, $\tan\angle AEC=\dfrac{1}{2}$, $\tan\angle ADC=1$, and the result says $\arctan\dfrac{1}{3}+\arctan\dfrac{1}{2}=45^\circ$.)
